\chapter{Pallor}

\section*{Definition}
Abnormal paleness of the skin and mucous membranes resulting from reduced hemoglobin or decreased cutaneous blood flow; clinically detectable when facial skin appears lighter than usual or conjunctival pallor is noted.

\section*{Pathophysiology}
Pallor arises from either decreased hemoglobin concentration (anemia reducing skin coloration) or vasoconstriction of cutaneous arterioles (reduced blood flow to skin surface). Anemia-induced pallor reflects reduced red cell mass, while vasoconstriction-mediated pallor reflects sympathetic activation or impaired perfusion in shock states.

\section*{Causes}
\begin{itemize}
  \item Anemia: Iron deficiency, vitamin B12/folate deficiency, thalassemia, sickle cell disease, acute blood loss, hemolysis, aplastic anemia, anemia of chronic disease
  \item Vasoconstriction: Shock (hypovolemic, cardiogenic, septic), cold exposure, syncope, fear or vasovagal episode, Raynaud phenomenon
  \item Endocrine: Hypopituitarism, hypothyroidism, Addison disease (darkening with pallor)
  \item Genetic: Albinism, vitiligo (depigmentation rather than true pallor)
\end{itemize}

\section*{When to See a Doctor}
See your doctor if:
\begin{itemize}
  \item Pallor with fatigue, dyspnea, tachycardia, or lightheadedness suggesting significant anemia
  \item Acute-onset pallor with hypotension, syncope, or signs of bleeding (hematemesis, melena, trauma)
  \item Progressive pallor with weight loss, easy bruising, or recurrent infections
  \item Pallor in infants or children with poor growth, pica, or developmental delay
\end{itemize}

\section*{Related Symptoms}
\begin{itemize}
  \item Fatigue
  \item Dyspnea on exertion
  \item Dizziness
  \item Cold intolerance
  \item Brittle nails (koilonychia in iron deficiency)
\end{itemize}
\textbf{Important:} Conjunctival pallor is the most reliable clinical sign of anemia; skin color alone can be misleading in dark-skinned individuals.

\section*{Self-Care / Home Management}
\begin{itemize}
  \item Eat an iron-rich diet if iron deficiency is suspected (red meat, leafy greens, legumes with vitamin C to enhance absorption).
  \item Avoid excessive caffeine or tannins (tea) with meals as they inhibit iron absorption.
  \item Stay warm and rest if pallor is due to cold exposure or vasovagal episode.
  \item Avoid self-medicating with iron supplements without blood tests confirming deficiency.
\end{itemize}

\section*{How Your Doctor Will Evaluate You}
\begin{itemize}
  \item Complete blood count (CBC) with red cell indices (MCV, MCH, RDW) to characterize anemia.
  \item Iron studies (ferritin, serum iron, TIBC), vitamin B12, folate levels based on indices.
  \item Peripheral blood smear to identify cell morphology (sickle cells, spherocytes, schistocytes).
  \item Reticulocyte count to distinguish hypo-proliferative from hemolytic or blood loss anemia.
\end{itemize}

\section*{Treatment Options}
\begin{itemize}
  \item Iron supplementation for iron deficiency anemia (oral ferrous sulfate 325 mg daily or every other day).
  \item Vitamin B12 injections or high-dose oral B12 for pernicious anemia; folate for folate deficiency.
  \item Blood transfusion for symptomatic anemia with hemodynamic instability or acute blood loss.
  \item Treat underlying cause: erythropoiesis-stimulating agents in CKD, disease-modifying therapy in hemolytic anemias.
\end{itemize}

\clearpage
